\documentclass[11pt]{article}
\usepackage[T1]{fontenc}
\usepackage{lmodern}
\usepackage{amsmath,amssymb,amsthm}
\usepackage[margin=1in]{geometry}
\usepackage[hidelinks]{hyperref}
\newtheorem{theorem}{Theorem}
\newtheorem{lemma}[theorem]{Lemma}
\newtheorem{proposition}[theorem]{Proposition}
\newtheorem{corollary}[theorem]{Corollary}
\theoremstyle{remark}
\newtheorem{remark}[theorem]{Remark}
\newcommand{\N}{\mathbb{N}}
\newcommand{\R}{\mathbb{R}}
\newcommand{\Z}{\mathbb{Z}}
\title{Positivity of a family of Stern polynomials\newline
and a counterexample to monotonicity}
\author{}
\date{13 October 2026}
\begin{document}
\maketitle
\begin{abstract}
For the standard Stern polynomials, we prove that
$B_{h_n}(t)>0$ for every real $t$ and every integer $n\geq0$, where
$h_n=\frac23(2^{2n}-1)(2^{2n+1}+1)+1$.
This proves the assertion of Conjecture~4.4 in the arXiv version of
Ulas's \emph{Strong arithmetic property of certain Stern polynomials}.
The proof derives from the defining recurrence a quadratic identity already
present in that source, and completes it to a sum of squares whose terms
cannot vanish simultaneously. We also compute $B_{183}$ and show that its
derivative at $-2/3$ is $-19/81$. An exact two-point certificate disproves
the final monotonicity clause of Conjecture~4.3 even on $(-1,0)$.
The general unique-real-root clause of that conjecture is not settled here.
A Lean~4 formalization certifies the all-parameter positivity theorem and
the exact monotonicity counterexample from the defining recurrence.
No claim of literature priority is made.
\end{abstract}

\section{Definitions and the statements under consideration}
Let $\N=\{0,1,2,\ldots\}$. The Stern polynomials $B_m\in\Z[t]$ are
defined by
\begin{equation}\label{eq:recurrence}
 B_0(t)=0,\qquad B_1(t)=1,\qquad
 B_{2m}(t)=tB_m(t),\qquad
 B_{2m+1}(t)=B_m(t)+B_{m+1}(t).
\end{equation}
The recurrence determines the polynomials when used for $m\geq1$;
both displayed identities remain valid for $m=0$.
This is the standard definition introduced by Klav\v{z}ar,
Milutinovi\'{c}, and Petr~\cite{KMP}, and is the definition in the body
of Ulas's paper~\cite[p.~2]{Ulas}. The odd-index formula in that paper's
abstract contains a different subscript; the present note uses the body
of the paper and \eqref{eq:recurrence}.

The sequence in Conjecture~4.4 of~\cite[p.~16]{Ulas} is
\begin{equation}\label{eq:h}
 h_n=\frac23(2^{2n}-1)(2^{2n+1}+1)+1\qquad(n\in\N).
\end{equation}
This is an integer sequence: if $K=4^n$ and $a=(K-1)/3$, then
$K\equiv1\pmod3$, $a\in\N$, and
\begin{equation}\label{eq:indices}
 K=3a+1,\qquad h_n=2a(2K+1)+1=4Ka+2a+1.
\end{equation}
In particular $h_0=1$, $h_1=19$, and $h_2=331$.
We will prove the stronger form of the conjectured absence of real roots:
\begin{equation}\label{eq:positive-goal}
 B_{h_n}(t)>0\qquad(n\in\N,\ t\in\R).
\end{equation}

For the monotonicity question, the family in~\cite[Theorem~2.6]{Ulas} is
\begin{equation}\label{eq:s}
 s_{0,n}=2^{2n+4}-9\cdot2^{n+1}-1\qquad(n\geq1).
\end{equation}
The last clause of Conjecture~4.3, numbered ``(3)'' again in the source,
asserts that $t\mapsto B_{s_{0,n}}(t)$ is increasing on $\R$.
The value $n=2$ lies in the range $n\geq2$ used for the first clause,
as well as in the positive-index range of the family. Thus it supplies
a counterexample under either of these index conventions.

\section{Two consequences of the defining recurrence}

\begin{lemma}[Concatenation identity]\label{lem:concat}
For integers $m,r,s\geq0$ with $s\leq2^r$, the following identity holds
in $\Z[t]$:
\begin{equation}\label{eq:concat}
 B_{m2^r+s}=B_{2^r-s}B_m+B_sB_{m+1}.
\end{equation}
All polynomials on the right have the same variable $t$.
\end{lemma}
\begin{proof}
We induct on $r$. If $r=0$, then $s=0$ or $s=1$;
the assertion follows from $B_0=0$ and $B_1=1$.
Suppose the result holds for $r$ and put $K=2^r$.
For an even remainder $s=2u$, where $0\leq u\leq K$, the induction
hypothesis and the even recurrence give
\begin{align*}
 B_{2mK+2u}
 &=t(B_{K-u}B_m+B_uB_{m+1})\\
 &=B_{2K-2u}B_m+B_{2u}B_{m+1}.
\end{align*}
For an odd remainder $s=2u+1\leq2K$, we have $0\leq u<K$.
Both $u$ and $u+1$ are admissible remainders at level $r$.
Consequently,
\begin{align*}
 B_{2mK+2u+1}
 &=B_{mK+u}+B_{mK+u+1}\\
 &=(B_{K-u}+B_{K-u-1})B_m+(B_u+B_{u+1})B_{m+1}\\
 &=B_{2K-2u-1}B_m+B_{2u+1}B_{m+1}.
\end{align*}
Here $K-u-1\geq0$, so the last use of the odd recurrence is legitimate.
This proves the induction step, including both endpoints of the range.
\end{proof}

\begin{lemma}[Adjacent values]\label{lem:adjacent}
If $t\in\R\setminus\{0\}$ and $m\in\N$, then
$(B_m(t),B_{m+1}(t))\neq(0,0)$. At $t=0$, every odd-indexed Stern
polynomial has value $1$.
\end{lemma}
\begin{proof}
Write $v_m=(B_m(t),B_{m+1}(t))^{\mathsf T}$. The recurrence gives
\begin{equation}\label{eq:pairs}
 v_{2m}=\begin{pmatrix}t&0\\1&1\end{pmatrix}v_m,
 \qquad
 v_{2m+1}=\begin{pmatrix}1&1\\0&t\end{pmatrix}v_m.
\end{equation}
For $t\neq0$, both matrices have nonzero determinant $t$ and hence
preserve the property of being a nonzero vector. Starting from
$v_0=(0,1)^{\mathsf T}$, strong induction on the index proves the
first assertion: for every positive index, its parent $\lfloor m/2\rfloor$
is smaller.

At $t=0$, the two transformations in \eqref{eq:pairs} take a pair
$(x,y)$ with $x+y=1$ to $(0,1)$ or $(1,0)$, respectively.
The initial pair has sum $1$, so the same induction shows that
$B_m(0)+B_{m+1}(0)=1$ for every $m$. The odd recurrence now gives
$B_{2m+1}(0)=1$.
\end{proof}

\section{The quadratic identity and strict positivity}

\begin{proposition}\label{prop:quadratic}
For $n\in\N$, let $a=(4^n-1)/3$, $A=B_a(t)$ and $C=B_{a+1}(t)$.
Then
\begin{equation}\label{eq:quadratic}
 B_{h_n}(t)=(t^2+t+1)A^2+(t+2)AC+C^2.
\end{equation}
\end{proposition}
\begin{proof}
Set $K=4^n=3a+1$. Since $2a+1\leq K\leq4K$ and
$4K=2^{2n+2}$, Lemma~\ref{lem:concat} applied to
$h_n=a(4K)+(2a+1)$ gives
\begin{equation}\label{eq:firstsplit}
 B_{h_n}=B_{4K-2a-1}A+B_{2a+1}C.
\end{equation}
The second coefficient is $B_{2a+1}=A+C$. To compute the first,
the odd recurrence gives
\[
 B_{4K-2a-1}=B_{2K-a-1}+B_{2K-a}.
\]
We have $K-a-1=2a$ and $K-a=2a+1$. Applying
Lemma~\ref{lem:concat} with $m=1$, $2^r=K$, and remainders
$2a$ and $2a+1$, which both lie in $[0,K]$, yields
\begin{align*}
 B_{2K-a-1}
 &=B_{K+2a}=B_{K-2a}B_1+B_{2a}B_2
   =C+t^2A,\\
 B_{2K-a}
 &=B_{K+2a+1}=B_{K-2a-1}B_1+B_{2a+1}B_2
   =A+t(A+C).
\end{align*}
Substitution in \eqref{eq:firstsplit} proves \eqref{eq:quadratic}.
The calculations also apply when $n=0$, for which $K=1$ and $a=0$.
\end{proof}

\begin{remark}[Attribution]
Ulas already gives identity \eqref{eq:quadratic} in the proof of
Theorem~3.1~\cite[p.~13]{Ulas}. No novelty is claimed for that identity.
The derivation above makes the
dependence on the original index sequence and defining recurrence
explicit. It does not use the subsequent displayed substitution or
initial values in that source.
\end{remark}

\begin{theorem}\label{thm:positivity}
For every $n\in\N$ and every $t\in\R$, $B_{h_n}(t)>0$.
In particular $B_{h_n}$ has no real root.
\end{theorem}
\begin{proof}
Completing the square in \eqref{eq:quadratic} gives
\begin{equation}\label{eq:sos}
 4B_{h_n}(t)=\bigl(2C+(t+2)A\bigr)^2+3t^2A^2.
\end{equation}
Both terms on the right are nonnegative. If $t\neq0$ and their
sum were zero, then $3t^2A^2=0$ would imply $A=0$. The other
square would then imply $C=0$. This contradicts
Lemma~\ref{lem:adjacent} for the adjacent indices $a,a+1$.
Thus $B_{h_n}(t)>0$ for $t\neq0$.

When $t=0$, formula \eqref{eq:indices} shows that $h_n$ is odd,
so Lemma~\ref{lem:adjacent} gives $B_{h_n}(0)=1$.
In particular the endpoint $n=0$ is covered: $h_0=1$ and
$B_{h_0}(t)=1$ for all $t$.
\end{proof}

\section{An exact failure of monotonicity}

\begin{proposition}\label{prop:counterexample}
For the sequence \eqref{eq:s}, $s_{0,2}=183$ and
\begin{equation}\label{eq:P}
 P(t):=B_{183}(t)=t^5+3t^4+7t^3+9t^2+5t+1.
\end{equation}
Moreover,
\begin{equation}\label{eq:derivative}
 P'(-2/3)=-19/81<0.
\end{equation}
The function $P$ is neither nondecreasing nor strictly increasing on
$(-1,0)$, and hence has neither property on $\R$.
\end{proposition}
\begin{proof}
First $s_{0,2}=2^8-9\cdot2^3-1=183$. For transparency, the
following intermediate values follow from \eqref{eq:recurrence}:
\begin{align*}
 B_5&=1+2t, & B_6&=t+t^2,\\
 B_{11}&=1+3t+t^2, & B_{12}&=t^2+t^3,\\
 B_{22}&=t+3t^2+t^3, & B_{23}&=1+3t+2t^2+t^3,\\
 B_{45}&=1+4t+5t^2+2t^3, & B_{46}&=t+3t^2+2t^3+t^4,\\
 B_{91}&=1+5t+8t^2+4t^3+t^4,
 & B_{92}&=t^2+3t^3+2t^4+t^5.
\end{align*}
Since $B_{183}=B_{91}+B_{92}$, these give \eqref{eq:P}.
Differentiating,
\[
 P'(t)=5t^4+12t^3+21t^2+18t+5,
\]
and exact substitution gives \eqref{eq:derivative}. Continuity of
$P'$ gives an interval around $-2/3$, contained in $(-1,0)$,
on which $P'<0$. The mean value theorem therefore gives strict
decrease on that interval.

There is also a certificate using only rational arithmetic:
\begin{gather*}
 -1<-\frac23<-\frac35<0,\qquad
 P(-2/3)=\frac{13}{243},\qquad P(-3/5)=\frac{122}{3125},\\
 P(-3/5)-P(-2/3)=-\frac{10979}{759375}<0.
\end{gather*}
These two points directly contradict both meanings of increasing
specified in the statement.
\end{proof}

\begin{proposition}\label{prop:root}
The polynomial $B_{183}$ has exactly one real root, namely $-1$,
and that root is simple.
\end{proposition}
\begin{proof}
Expansion verifies the identity
\[
 P(t)=(t+1)\left([t(t+1)]^2+(2t+1)^2\right).
\]
The first square can vanish only at $0$ or $-1$, whereas the second
can vanish only at $-1/2$. Their sum is therefore strictly positive
on $\R$. This proves that $P(t)=0$ precisely when $t=-1$.
The factor in parentheses has value $1$ at $-1$, so the root is simple.
\end{proof}

Proposition~\ref{prop:root} prevents a possible misinterpretation of the
counterexample: the failure of monotonicity does not refute the
unique-real-root clause of Conjecture~4.3. No assertion for the general
parameters of that clause is proved here. Finite root counts do not
establish that general assertion.

\section{Formal verification and reproducibility}
The accompanying project has been checked using Lean~4.19.0 and
mathlib~4.19.0. It defines a real-valued function $B$ recursively and
proves the standard recurrence. The theorem
\path{Stern.B_integer_polynomial} constructs, for every index,
an integer polynomial whose real evaluation is $B$.
The formal identities are pointwise equalities of real-valued functions;
the project does not separately construct a canonical sequence of objects
in $\Z[X]$ and prove coefficient equalities between those objects.

The source index is defined using the explicit formula
\eqref{eq:h}, with natural-number division.
\path{Stern.h_source_exact} proves that its rational value is
exactly \eqref{eq:h}; thus integer truncation does not alter the sequence.
The recurrence leads to the concatenation and adjacent-value lemmas,
then to \path{Stern.h_quadratic} and
\path{Stern.h_sum_of_squares}.
The principal theorem is
\[
\texttt{Stern.h\_positive}:
\qquad \forall n\in\N\;\forall t\in\R,\quad 0<B_{h_n}(t).
\]
Its zero-parameter case and absence-of-real-roots corollary are also
checked as \path{Stern.h_at_zero} and
\path{Stern.h_no_real_root}.

The counterexample file checks the recurrence computation
\path{Stern.B183_explicit}, the derivative formula, and
\path{Stern.B183_deriv_neg_two_thirds}.
It connects the negative derivative directly to
\path{Stern.B183_not_monotoneOn} on $(-1,0)$ by showing that a
nondecreasing function's right-hand slopes cannot have a negative limit.
Failure of strict increase and of global increase are also checked,
as are the original-family specialization $s_{0,2}=183$ and the exact
two-point certificate \path{Stern.B183_two_point_certificate}.
Finally, \path{Stern.B183_zero_iff} proves that a real zero is
exactly $-1$, using the verified factorization and positive quotient.
The theorem \path{Stern.B183_deriv_neg_one} also checks $P'(-1)=1$.
The elementary simplicity argument in Proposition~\ref{prop:root}
is not separately encoded as a polynomial-root multiplicity theorem.

All five library modules, their aggregate module, and the axiom audit
compiled successfully. No admitted proof or additional axiom is used.
The printed axiom dependencies are subsets of the standard Lean
dependencies \texttt{propext}, \texttt{Classical.choice}, and
\texttt{Quot.sound}. Exact theorem statements, the tested commands,
and their outputs are recorded in \texttt{STATUS.md}, \texttt{BUILD.log},
and \texttt{AXIOMS.log}. The local build used the installed toolchain
and existing mathlib dependencies; a fresh network download of those
dependencies was not tested.
The supplementary Python checker uses exact arithmetic for finite
checks, including $0\leq n\leq20$. Neither the natural-language proof
nor the formal all-parameter theorem depends on those finite tests.

\section{Sources, priority, and assistance}
The conjecture statements and the quadratic identity were checked against
the 17-page arXiv version~\cite{Ulas}. The journal publication metadata
are supplied in the bibliography, but access to its full text was
restricted; page-by-page agreement with the journal version is not
asserted. The accessible paper by Altizio~\cite{AltizioPaper} was also
examined. Its stated results on zero-free complex regions and
prime-index irreducibility do not by themselves establish the results
proved here. The official record for Altizio's 2025 doctoral
thesis~\cite{AltizioThesis} was available, but the thesis full text
could not be accessed. Its coverage of the present questions is
therefore undetermined. The literature search is not exhaustive, and
this note does not assert that these conclusions are the first published
resolutions. This work was prepared with AI assistance; no claim of
human expert peer review is made.

\begin{thebibliography}{9}
\bibitem{KMP}
S.~Klav\v{z}ar, U.~Milutinovi\'{c}, and C.~Petr,
\emph{Stern polynomials}, Advances in Applied Mathematics
\textbf{39} (2007), 86--95.

\bibitem{Ulas}
M.~Ulas, \emph{Strong arithmetic property of certain Stern polynomials},
arXiv:1909.10844v1 (2019),
\url{https://arxiv.org/abs/1909.10844v1}.
Published in Publicationes Mathematicae Debrecen
\textbf{96} (3--4) (2020), 401--422,
\href{https://doi.org/10.5486/PMD.2020.8683}{doi:10.5486/PMD.2020.8683}.
Page references in this note are to the arXiv version.

\bibitem{AltizioPaper}
D.~Altizio, \emph{Zeros of Stern polynomials in the complex plane},
arXiv:2511.03847v1 (2025),
\url{https://arxiv.org/abs/2511.03847v1}.

\bibitem{AltizioThesis}
D.~Altizio, \emph{On Stern polynomials and their zeroes},
Ph.D. thesis, University of Illinois Urbana-Champaign, 2025.
Official record: \url{https://www.ideals.illinois.edu/items/136452}.
Metadata consulted; full text unavailable for this note.
\end{thebibliography}
\end{document}
